% SPDX-License-Identifier: LPPL-1.3c
% Copyright (C) 2026 Christophe Jorssen
% Author and Current Maintainer: Christophe Jorssen <christophe.jorssen@gmail.com>
% LPPL maintenance status: maintained
% This file is part of luafemm. See LICENSE and LICENSES.md for its terms.
\section{Getting started}
\label{sec:start}
\begin{tikzlibrary}{femm}
  This library gives evaluated TikZ paths a physical meaning. It supplies
  a planar magnetostatic model, a constrained triangular mesh, material
  choices, oriented field lines and measurement paths. Computation runs
  inside LuaTeX; no native solver or shell escape is required.
\end{tikzlibrary}

\subsection{Requirements and installation}
Use a LuaTeX engine with Lua 5.3 or later, PGF/TikZ, and optionally PGFPlots
for axes. The release was tested with LuaTeX/LuaHBTeX 1.24.0, PGF 3.1.12,
PGFPlots 1.18.3 and ConTeXt MkIV 2026.09.10. It does not support pdfTeX,
XeTeX, or ConTeXt MkII. ConTeXt LMTX is not part of the release test matrix.

Copy the distribution's |tex| tree into a TEXMF tree; keep all Lua and
generic TeX files together under |tex/generic/luafemm|. The LaTeX wrapper
belongs in |tex/latex/luafemm| and the ConTeXt wrapper in
|tex/context/third/luafemm|. Refresh the filename database if your TeX
distribution requires it. Keep the supplied licence and provenance files
with redistributed copies. From a source checkout, |make test|,
|make examples| and |make manual| set up the search paths for you.

\subsection{LaTeX}
Either load the convenience package |luafemm| or load the libraries directly:
\begin{codeexample}[code only]
\documentclass{article}
\usepackage{tikz,pgfplots}
\usetikzlibrary{femm}
\usepgfplotslibrary{femm}
\pgfplotsset{compat=1.18}
\begin{document}
% Place a tikzpicture here.
\end{document}
\end{codeexample}
Compile with |lualatex --no-shell-escape document.tex|.
The package wrapper only loads TikZ and the generic library; PGFPlots
remains optional.

\subsection{Plain LuaTeX}
\begin{codeexample}[code only]
\input tikz.tex
\input pgfplots.tex
\usetikzlibrary{femm}
\usepgfplotslibrary{femm}
\pgfplotsset{compat=1.18}
\tikzpicture[femm/problem]
  % Physical paths and field pictures.
\endtikzpicture
\bye
\end{codeexample}
Use |luatex --no-shell-escape document.tex|. The macros are identical to
the LaTeX ones; |\tikzpicture|/|\endtikzpicture| and |\axis|/|\endaxis|
replace the environment syntax.

\subsection{ConTeXt MkIV}
\begin{codeexample}[code only]
\usemodule[luafemm]
\usemodule[pgfplots]
\usepgfplotslibrary{femm}
\pgfplotsset{compat=1.18}
\starttext
\starttikzpicture[femm/problem]
  % Physical paths and field pictures.
\stoptikzpicture
\stoptext
\end{codeexample}
Compile with |context --luatex document.tex|, explicitly selecting MkIV.
Use |\startaxis|/|\stopaxis| inside a |\starttikzpicture| for a graph.
Loading |\usemodule[tikz]| followed by |\usetikzlibrary[femm]| is an
alternative to the wrapper. Repeated loads are harmless.

\subsection{A complete small model}
The following example is executed while this manual is built. Its red
rectangle is a permanent magnet pointing up. Every path declared as a
region is also used by the mesher.
\begin{codeexample}[width=5cm]
\begin{tikzpicture}[
  femm/problem={xmin=-15,xmax=15,
    ymin=-15,ymax=15,mesh size=2}]
  \filldraw[femm/region={
    material=cast-alnico-5-lng37,
    magnetization angle=90},
    fill=red!20]
    (-3,-5) rectangle (3,5);
  \pic[femm/lines=13] {femm field};
\end{tikzpicture}
\end{codeexample}

The workflow is: define any custom materials, create a problem in the
picture options, declare all physical regions, and then request a mesh,
field or profile plot. The first such request computes what it needs.
Once meshed, a model cannot accept additional materials or regions.
Measurement paths may be registered before or after solving.

\subsection{Notation, scope and stability}
Options have the usual PGF |key=value| syntax. Most public keys start with
|/tikz/femm|; lists introduced by |femm/problem|, |femm/region| and
|femm/profile| have their own documented subkeys. Settings follow TeX
grouping and may be inherited from styles or scopes. A new named profile
is a Lua object with document-wide lifetime; this is deliberately different
from a local TeX option.

The defaults in this manual are initial values unless described as a reset
performed by a command. A ``number'' means a finite decimal or scientific
literal without a TeX unit. Only keys explicitly documented as accepting
PGF expressions evaluate such expressions. Graphic dimensions such as
|arrow length| follow ordinary TikZ dimension syntax.

The release-candidate public interface comprises the keys and commands
documented here. Macros containing |\luafemm@|, capture postactions, and
the internal mesh layout are not extension interfaces. The Lua API is
documented separately in Section~\ref{sec:development}; its internal
tables should not be edited after meshing.
